An app that spends for a user today either asks the user to approve every call
or holds a key that can do anything. A Soroban account on Stellar can be a
contract that decides each authorization itself (\S\ref{sec:background}), so a
middle ground exists: a \emph{grant} lets an app or agent act for the account
within limits the account itself enforces. The bound has to hold across the
nested calls of one transaction, across repeated calls over time, and whichever
installed rule the transaction's submitter selects.

Our running example, the shop grant (Fig.~\ref{fig:fig-shop-grant}), lets an
app or agent send XLM from the account to the shop and to no one else, at most
\fact{shop-cap-xlm}~XLM in any \fact{shop-period-ledgers} ledgers
(\fact{shop-period-minutes}). A grant cannot say everything: it has no bound on
one call's amount (the cap bounds the window), no \emph{or} and no \emph{not},
no reads of other contracts, and it checks each call alone.

\perchfigure{fig-shop-grant}

We present Perch, a Soroban smart account that offers safe, bounded delegation
for apps and agents: simple, proved, and testable end to end in soroban-sim.
\emph{Safe} and \emph{bounded} mean three properties, not an audit: each call of
an accepted request is allowed by an installed rule, no window of a capped rule holds more than its
limit, and no grant allows a call to the account or its adopted recovery
contracts (Theorems~\ref{stmt:acceptance-implies-permission},
\ref{stmt:cap-window-safety} and~\ref{stmt:grant-target-containment}).
\emph{Simple} means one document, one install call
(\code{apply\_scope}) and one decision point, at the cost of fewer conditions.
\emph{Proved} applies to the Lean translation of the Rust decision function,
not the deployed Wasm or the tested account code around it, the \emph{shell}. \emph{End to end} covers both designs' contract Wasm and
the token in one local host, not the wallet, RPC, passkeys or fees. A completed
recovery moves one counter, and every earlier grant stops authorizing.

\paragraph{Status.} The design is an unmerged pull-request stack ending at
\fact{perch-pr-stack-head}, previewed on testnet through
\fact{wallet-name}~\selfcite{nido-preview}, and not audited; we do not recommend it for real
funds before an audit. An earlier design is on mainnet~\selfcite{perch-deployments}. \ifdefined\anonymous\else
soroban-sim~\cite{soroban-sim} is published by The Aha Company, where the
authors work.\fi

\paragraph{Contributions.} (1)~An account design with one decision point and
account-held spending windows (\S\S\ref{sec:grants}--\ref{sec:authorization});
(2)~Lean theorems about the translation of its decision function
(\S\ref{sec:proofs}); (3)~a dual soroban-sim simulation of Perch and an
OpenZeppelin account with oz-policy-builder from the same contract bytes
(\S\ref{sec:simulation}); and (4)~a source-cited comparison, including where
oz-policy-builder is ahead (\S\ref{sec:compare}). We ask whether the account
enforces grants and caps as designed (\textbf{RQ1}, \S\S\ref{sec:simulation}
and~\ref{sec:evaluation}), what is proved, tested or trusted (\textbf{RQ2},
\S\ref{sec:proofs}), and how the designs differ on one task (\textbf{RQ3},
\S\S\ref{sec:compare}--\ref{sec:simulation}).
Companions cover the comparison, \perchdoc{oz-policy-builder}; all theorems,
\perchdoc{verification}; and recovery and agents, \perchdoc{vision}.
